\Article{7}
Nach \textsc{d'Alembert} veröffentlichte der ausgezeichnete \textsc{Euler} seine Untersuchungen über denselben Gegenstand in \textit{Recherches sur les racines imaginaires des équations, Hist. de l'Acad. de Berlin A.} 1749, S. 223 ff. Er gab dort zwei Methoden an; der wesentliche Inhalt der ersten ist folgender.

Zunächst unternimmt der ausgezeichnete Euler den Beweis, dass, wenn $m$ irgendeine Potenz der Zahl 2 bezeichnet, die Funktion $x^{2m}+Bx^{2m-2}+Cx^{2m-3}+\text{usw.}+M=X$ (in der der Koeffizient des zweiten Gliedes $=0$ ist) stets in zwei reelle Faktoren zerlegt werden kann, in denen $x$ bis zum Grad $m$ aufsteigt. Dazu nimmt er zwei Faktoren an,
\[
x^m-ux^{m-1}+\alpha x^{m-2}+\beta x^{m-3}+\text{usw.},\;\text{und}\; x^m+ux^{m-1}+\lambda x^{m-2}+\mu x^{m-3}\ \text{usw.}
\]
wobei die Koeffizienten $u,\alpha,\beta$ usw., $\lambda,\mu$ usw. noch unbekannt sind, und setzt ihr Produkt der Funktion $X$ gleich. Der Vergleich der Koeffizienten liefert dann $2m-1$ Gleichungen. Offenbar muss nur bewiesen werden, dass man den Unbekannten $u,\alpha,\beta$ usw., $\lambda,\mu$ usw. (deren Anzahl ebenfalls $2m-1$ ist) solche reellen Werte geben kann, die jene Gleichungen erfüllen. Euler behauptet nun: Betrachtet man zunächst $u$ als bekannt, sodass die Anzahl der Unbekannten um eins kleiner ist als die der Gleichungen, so lassen sich durch geeignete Verbindung dieser Gleichungen nach bekannten algebraischen Methoden alle $\alpha,\beta$ usw., $\lambda,\mu$ usw. rational und ohne irgendein Wurzelziehen durch $u$ und die Koeffizienten $B,C$ usw. bestimmen; sie erhalten also reelle Werte, sobald $u$ reell wird. Außerdem kann man alle $\alpha,\beta$ usw., $\lambda,\mu$ usw. eliminieren, sodass eine Gleichung $U=0$ entsteht, in der $U$ ein Polynom allein in $u$ und den bekannten Koeffizienten ist. Diese Gleichung selbst nach der gewöhnlichen Eliminationsmethode zu entwickeln, wäre eine ungeheure Arbeit, wenn die vorgelegte Gleichung $X=0$ von etwas\SourceLeaf{017} höherem Grad ist, und für einen unbestimmten Grad ganz unmöglich (nach Eulers eigenem Urteil, S. 239). Hier genügt es jedoch, eine Eigenschaft dieser Gleichung zu kennen: Das konstante Glied in $U$ (das die Unbekannte $u$ nicht enthält) ist notwendig negativ. Daraus folgt bekanntlich, dass die Gleichung wenigstens eine reelle Wurzel hat oder dass sich $u$ und somit auch $\alpha,\beta$ usw., $\lambda,\mu$ usw. auf wenigstens eine Weise reell bestimmen lassen. Jene Eigenschaft lässt sich durch die folgenden Überlegungen bestätigen. Da $x^m-ux^{m-1}+\alpha x^{m-2}+\text{usw.}$ ein Faktor der Funktion $X$ sein soll, muss $u$ die Summe von $m$ Wurzeln der Gleichung $X=0$ sein und daher ebenso viele Werte haben, wie man auf verschiedene Weise aus $2m$ Wurzeln $m$ auswählen kann, also nach den Grundsätzen der Kombinatorik $\frac{2m\mathbin{.}\overline{2m-1}\mathbin{.}\overline{2m-2}.\,.\,.\,.\,.\overline{m+1}}{1\mathbin{.}2\mathbin{.}3\ldots m}$ Werte. Diese Zahl ist stets das Doppelte einer ungeraden Zahl (den nicht schwierigen Beweis lasse ich weg). Setzt man sie also $=2k$, so ist ihre Hälfte $k$ ungerade; die Gleichung $U=0$ hat aber den Grad $2k$. Da nun in der Gleichung $X=0$ das zweite Glied fehlt, ist die Summe aller $2m$ Wurzeln 0. Ist also die Summe irgendwelcher $m$ Wurzeln $+p$, so ist die Summe der übrigen $-p$; das heißt: Gehört $+p$ zu den Werten von $u$, so gehört auch $-p$ dazu. Hieraus schließt Euler, dass $U$ das Produkt aus $k$ quadratischen Faktoren wie $uu-pp$, $uu-qq$, $uu-rr$ usw. ist, wobei $+p$, $-p$, $+q$, $-q$ usw. alle $2k$ Wurzeln der Gleichung $U=0$ bezeichnen. Wegen der ungeraden Anzahl dieser Faktoren ist daher das konstante Glied in $U$ das Quadrat des Produkts $pqr$ usw. mit negativem Vorzeichen. Das Produkt $pqr$ usw. lässt sich aber stets aus den Koeffizienten $B,C$ usw. rational bestimmen und ist deshalb notwendig eine reelle Größe. Sein Quadrat mit negativem Vorzeichen ist somit gewiss eine negative Größe. Q. E. D.

Da diese beiden reellen Faktoren von $X$ den Grad $m$ haben und $m$ eine Potenz der Zahl 2 ist, lässt sich jeder auf dieselbe Weise erneut in zwei reelle Faktoren vom Grad $\frac12m$ zerlegen. Durch wiederholte Halbierung der Zahl $m$ muss man schließlich zu zwei gelangen; daher erhält man durch Fortsetzung des Verfahrens die Funktion $X$ offenbar in reelle quadratische Faktoren zerlegt.

Wird dagegen eine Funktion vorgelegt, in der das zweite Glied nicht fehlt, etwa $x^{2m}+Ax^{2m-1}+Bx^{2m-2}+\text{usw.}+M$, wobei $2m$ weiterhin eine Zweierpotenz bezeichnet, so geht sie durch die Substitution $x=y-\frac{A}{2m}$ in eine ähnliche Funktion ohne zweites Glied über. Daraus folgt leicht, dass auch diese Funktion in reelle quadratische Faktoren zerlegt werden kann.

Ist schließlich eine Funktion vom Grad $n$ vorgelegt, wobei $n$ eine Zahl bezeichnet, die keine
\SourceLeaf{018}
Zweierpotenz ist, so setze man die nächstgrößere Zweierpotenz über $n$ gleich $=2m$ und multipliziere die vorgelegte Funktion mit beliebigen $2m-n$ reellen linearen Faktoren. Aus der Zerlegbarkeit des Produkts in reelle quadratische Faktoren ergibt sich ohne Schwierigkeit, dass auch die vorgelegte Funktion in reelle Faktoren zweiten oder ersten Grades zerlegbar sein muss.
